% Extracción quirúrgica del propietario V2 05d: líneas 286--512.
% El resto permanece en este archivo como procedencia, pero no se publica.
\iffalse
% Macroestructura de coherencia estrictamente interna del continuo discreto.
% No usa geometrías, clases ni conclusiones clásicas.

\chapter{Internal macrostructure of cells, coherence and completeness}
\label{chap:pdfv2-coh-interna}

\begin{statusbox}
The sole input of this chapter is the discrete structure of the continuum. The
branch \(\mathrm{coh}\) is constructed on the rings
\(A_N=\mathbb Z/3^N\mathbb Z\), the APP cells, their histories, and their
terminals. Its nonablative chain is
\[
 \mathcal C_{\rm cont}^{\rm disc}
 \supset\mathcal C_{\rm cont}^{\rm coh}
 \xrightarrow{\operatorname{Res}_{\rm coh}}
 \mathcal G_{\rm coh}
 \xrightarrow{\operatorname{pr}_{X,{\rm coh}}}
 X_{\chi,{\rm coh}}
 \xrightarrow{\mathcal A_{\rm coh}}
 \mathfrak M_{\rm coh}
 \xrightarrow{\mathcal P_{\rm coh}}
 \mathcal C_{\rm coh}^{\rm HMT}.
\]
No classical object appears as an input or as an output. Each image literally preserves its generating history.
\end{statusbox}

\section{Typed domain and index separation}
\label{sec:pdfv2-coh-domain}

To avoid identifying depth with precision, the directed set is used.
\begin{equation}\label{eq:pdfv2-coh-lambda}
 \Lambda:=\mathbb N_{\rm prof}\times\mathbb N_{\rm prec},
 \qquad
 \lambda=(d,N),
 \qquad
 A_N:=\mathbb Z/3^N\mathbb Z.
\end{equation}
If \(\lambda\leq\lambda'=(d',N')\), the transition
\begin{equation}\label{eq:pdfv2-coh-rho-lambda}
 \rho_{\lambda',\lambda}:
 \mathcal C_{{\rm cont},\lambda'}^{\rm disc}
 \longrightarrow
 \mathcal C_{{\rm cont},\lambda}^{\rm disc}
\end{equation}
truncates the history after depth \(d\) and reduces coefficients by
\(A_{N'}\twoheadrightarrow A_N\), preserving sheet, orientation, residue,
quotient, carry, memory, boundary, path, survival, and terminal.

The substructure \(\mathcal C_{\rm cont}^{\rm coh}\) consists of the histories
\(x\in\mathcal C_{\rm cont}^{\rm disc}\) that satisfy, at every level:
\begin{enumerate}
 \item the prefix of \(x\) carries the oriented and based cells defined in
       the following section;
 \item each refinement induce a map of the complex of the celda;
 \item truncation, coefficient reduction, orientation, and nonadic
       extension commute;
 \item the surviving-extension multisections are nonempty at every
       depth;
 \item the nonadic return preserves the terminal difference between the history and its continuation.
\end{enumerate}
These conditions define the domain of the branch. They are not attributed to an
arbitrary history without verifying them.

\section{APP cell over \texorpdfstring{\(A_N\)}{AN}}
\label{sec:pdfv2-coh-cell}

An oriented cell is
\[
 \nu=(h_\nu,I_\nu,J_\nu,\epsilon_\nu),
\]
where \(h_\nu\) is a finite history, and \(I_\nu,J_\nu\) are nonempty
ordered sets of ports,
\[
 a_\nu:=|I_\nu|\geq1,
 \qquad
 b_\nu:=|J_\nu|\geq1,
\]
and \(\epsilon_\nu\in\{+1,-1\}\) is its orientation.

We define
\begin{equation}\label{eq:pdfv2-coh-kappa}
 \begin{aligned}
 \kappa_{\nu,N}:
 A_N^{I_\nu}\oplus A_N^{J_\nu}
 &\longrightarrow A_N^{I_\nu\times J_\nu},\\
 \kappa_{\nu,N}(u,v)_{ij}&:=u_i-v_j.
 \end{aligned}
\end{equation}
Once the base ports \(i_0\in I_\nu\), \(j_0\in J_\nu\) have been chosen, let
\begin{equation}\label{eq:pdfv2-coh-delta}
 \begin{aligned}
 \delta_{\nu,N}:A_N^{I_\nu\times J_\nu}
 &\longrightarrow
 A_N^{(I_\nu\setminus i_0)\times(J_\nu\setminus j_0)},\\
 \delta_{\nu,N}(w)_{ij}
 &:={}
 w_{ij}-w_{i j_0}-w_{i_0j}+w_{i_0j_0}.
 \end{aligned}
\end{equation}

\begin{proposition}[Exact sequence of the cell]
\label{prop:pdfv2-coh-cell-exact}
For every celda \(\nu\) and every precision \(N\), the sucesion
\begin{equation}\label{eq:pdfv2-coh-exact-sequence}
 0\longrightarrow A_N
 \xrightarrow{\ \Delta\ }
 A_N^{I_\nu}\oplus A_N^{J_\nu}
 \xrightarrow{\ \kappa_{\nu,N}\ }
 A_N^{I_\nu\times J_\nu}
 \xrightarrow{\ \delta_{\nu,N}\ }
 A_N^{(a_\nu-1)(b_\nu-1)}
 \longrightarrow0,
\end{equation}
where
\[
 \Delta(t)=\bigl((t)_{i\in I_\nu},(t)_{j\in J_\nu}\bigr),
\]
is exact and split over \(A_N\).
\end{proposition}

\begin{proof}
If \(\kappa_{\nu,N}(u,v)=0\), then \(u_i=v_j\) for every \(i,j\), so \((u,v)\) is diagonal. Thus,
\(\ker\kappa_{\nu,N}=\operatorname{im}\Delta\). Direct substitution gives
\(\delta_{\nu,N}\kappa_{\nu,N}=0\).

If \(\delta_{\nu,N}(w)=0\), set
\[
 u_i=w_{i j_0},
 \qquad
 v_j=w_{i_0j_0}-w_{i_0j}.
\]
Then
\[
 u_i-v_j=w_{i j_0}+w_{i_0j}-w_{i_0j_0}=w_{ij},
\]
so that \(\ker\delta_{\nu,N}=\operatorname{im}\kappa_{\nu,N}\). To prove the surjectivity of \(\delta_{\nu,N}\), the base
row and column are set to zero and the desired coordinates are placed
arbitrarily in the remaining positions. Those same formulas provide
splittings.
\end{proof}

The complex cell, in degrees \(1\to0\), is the homological cone of
\(\kappa_{\nu,N}\) with the convention that places its source in degree \(1\):
\begin{equation}\label{eq:pdfv2-coh-cell-complex}
 C_{\nu,N}:=\operatorname{Cone}(\kappa_{\nu,N})
 :=\left[
 A_N^{I_\nu}\oplus A_N^{J_\nu}
 \xrightarrow{\ \kappa_{\nu,N}\ }
 A_N^{I_\nu\times J_\nu}
 \right].
\end{equation}
The preceding proposition gives
\begin{equation}\label{eq:pdfv2-coh-cell-homology}
 H_1(C_{\nu,N})\simeq A_N,
 \qquad
 H_0(C_{\nu,N})\simeq A_N^{(a_\nu-1)(b_\nu-1)}.
\end{equation}
The cell defect is thus constructed, not added as a cardinal without
map.

\section{Orientation and carry of the cell}
\label{sec:pdfv2-coh-orientation-carry}

Let
\[
 \tau_1(u,v)=(v,u),
 \qquad
 P(w)_{ji}=w_{ij}.
\]
Oriented exchange is the complex map
\begin{equation}\label{eq:pdfv2-coh-theta}
 \Theta_{\nu,N}:=(\tau_1,-P):
 C_{a_\nu,b_\nu,N}\longrightarrow C_{b_\nu,a_\nu,N},
\end{equation}
because
\begin{equation}\label{eq:pdfv2-coh-orientation}
 \kappa_{b_\nu,a_\nu,N}\tau_1
 =-P\kappa_{a_\nu,b_\nu,N}.
\end{equation}
Moreover,
\(\Theta_{\nu^{\rm op},N}\Theta_{\nu,N}=I\). The opposite route is not
identified with the original: it changes the degree-zero sign and preserves the act
of inversion.

For \(r\in A_N\), let
\(\langle r\rangle_N\in\{0,\ldots,3^N-1\}\) be its canonical representative. The
local borrow/carry is
\begin{equation}\label{eq:pdfv2-coh-beta}
 \beta_N(r,s):=
 \frac{\langle r-s\rangle_N-\langle r\rangle_N+\langle s\rangle_N}
 {3^N}\in\{0,1\}.
\end{equation}
Thus,
\begin{equation}\label{eq:pdfv2-coh-carry-cell}
 \langle u_i-v_j\rangle_N
 =\langle u_i\rangle_N-\langle v_j\rangle_N
  +3^N\beta_N(u_i,v_j).
\end{equation}
The matrix
\[
 \operatorname{Carr}_{\nu,N}(u,v)
 :=\bigl(\beta_N(u_i,v_j)\bigr)_{ij}
\]
is preserved alongside \(\kappa_{\nu,N}(u,v)\): residue and carry are two
distinct coordinates.

When port sizes follow from the two APP leaves, the admissibility
of the cell includes the integer identity
\begin{equation}\label{eq:pdfv2-coh-app-defect}
 (a_\nu-1)(b_\nu-1)
 =\rho_{\Pi,\nu}-\rho_{\Sigma,\nu}+1
  +9(c_{\Pi,\nu}-c_{\Sigma,\nu}).
\end{equation}
The equality is required only for a cell with that genealogy; it is not asserted
for integers devoid of APP leaves.

If \(L_N(u)\) is the lift whole canonica of the matrix of a flecha
\(u\), it define the carry of composition by
\begin{equation}\label{eq:pdfv2-coh-carry-composition}
 L_N(v)L_N(u)=L_N(vu)+3^N c_N(v,u).
\end{equation}
The associativity of three arrows gives the exact cocycle.
\begin{equation}\label{eq:pdfv2-coh-carry-cocycle}
 c_N(w,vu)+L_N(w)c_N(v,u)
 =c_N(wv,u)+c_N(w,v)L_N(u).
\end{equation}

\section{Category of histories and cellular functor}
\label{sec:pdfv2-coh-category}

For a prefix \(x_{\leq d}\), let \(J_d(x)\) be the finite category whose
objects are all its cells \(\nu\). Its arrows are generated by:
\begin{enumerate}
 \item refinements \(u:\nu\to\nu'\) with surjective maps of
       ports
       \[
        \alpha_u:I_{\nu'}\twoheadrightarrow I_\nu,
        \qquad
        \beta_u:J_{\nu'}\twoheadrightarrow J_\nu;
       \]
 \item the exchanges of orientation \(\Theta_\nu\);
 \item the nonadic extensions \(\Gamma_9\nu\);
 \item identity and composition, subject to the boundary,
       path, orientation, and carry relations of the history.
\end{enumerate}
Finiteness refers to the prefix; the system of all prefixes is not declared finite.

The cellular functor is
\begin{equation}\label{eq:pdfv2-coh-gamma-functor}
 \Gamma_{d,N}:J_d(x)\longrightarrow
 \operatorname{Ch}^{[0,1]}(A_N),
 \qquad
 \Gamma_{d,N}(\nu)=C_{\nu,N}.
\end{equation}
For a refinement arrow \(u:\nu\to\nu'\), it acts by means of
\begin{align}
 \Gamma_{d,N}(u)_1(p,q)
 &:=(p\circ\alpha_u,q\circ\beta_u),
 \label{eq:pdfv2-coh-gamma-one}\\
 \Gamma_{d,N}(u)_0(w)
 &:=w\circ(\alpha_u\times\beta_u).
 \label{eq:pdfv2-coh-gamma-zero}
\end{align}
It is verified coordinate by coordinate.
\begin{equation}\label{eq:pdfv2-coh-gamma-chain-map}
 \kappa_{\nu',N}\Gamma_{d,N}(u)_1
 =\Gamma_{d,N}(u)_0\kappa_{\nu,N}.
\end{equation}
On an inversion of orientation it acts by
\(\Theta_{\nu,N}\), and on compositions it acts by composition, separately preserving
\(c_N(v,u)\).

The reduction \(r_{N+1,N}:A_{N+1}\twoheadrightarrow A_N\) satisfies
\begin{equation}\label{eq:pdfv2-coh-kappa-reduction}
 r_{N+1,N}^{I\times J}\kappa_{\nu,N+1}
 =\kappa_{\nu,N}
  \bigl(r_{N+1,N}^{I}\oplus r_{N+1,N}^{J}\bigr).
\end{equation}
Therefore, the cell diagram is naturally simultaneous in history and
precision.

\fi
\section{Module of paths, cellular totalization and nonadic terminal}
\label{sec:pdfv2-coh-route-module}

Let \(\operatorname{Term}_d(x)\) be the conjunto of terminal of memory of the
prefix. For \(\nu\in J_d(x)\), definase
\begin{equation}\label{eq:pdfv2-coh-w-module}
 W_{d,N}(\nu):=
 A_N\bigl[
 \operatorname{Path}_{J_d(x)}(\nu,\operatorname{Term}_d(x))
 \bigr].
\end{equation}
It is the free module on all terminal routes; it selects none. If
\(u:\nu\to\nu'\), the contravariant action is
\begin{equation}\label{eq:pdfv2-coh-w-action}
 W_{d,N}(u):W_{d,N}(\nu')\longrightarrow W_{d,N}(\nu),
 \qquad
 [\lambda]\longmapsto[\lambda\circ u].
\end{equation}
Each generator preserves orientation, carry, boundary, route, and memory.

\section{Bilateral bar complex and totalization of the bicolumn}
\label{sec:pdfv2-coh-bar}

For \(q\geq0\) and \(r\in\{0,1\}\), let
\begin{equation}\label{eq:pdfv2-coh-bar-bigraded}
 B_{q,r}^{d,N}(x):=
 \bigoplus_{\nu_0\xrightarrow{u_1}\cdots\xrightarrow{u_q}\nu_q}
 W_{d,N}(\nu_q)\otimes_{A_N}\Gamma_{d,N}(\nu_0)_r.
\end{equation}
The bar differential is defined without abbreviation by
\begin{align}
 d_{\rm bar}(w;u_1,\ldots,u_q;c)
 :={}&(w;u_2,\ldots,u_q;\Gamma(u_1)c)
 \nonumber\\
 &+\sum_{i=1}^{q-1}(-1)^i
 (w;u_1,\ldots,u_{i+1}u_i,\ldots,u_q;c)
 \nonumber\\
 &+(-1)^q
 (W(u_q)w;u_1,\ldots,u_{q-1};c).
 \label{eq:pdfv2-coh-bar-differential}
\end{align}
For \(q=1\), this is exactly the relation
\begin{equation}\label{eq:pdfv2-coh-coend-relation}
 d_{\rm bar}(w;u;c)
 =w\otimes\Gamma(u)c-W(u)w\otimes c.
\end{equation}
A transported incidence is thereby identified with the same incidence written
in the refined chart; the relation is not concealed beneath the name of the
totalizer.

The total complex is
\begin{equation}\label{eq:pdfv2-coh-total-complex}
 K_{d,N}(x):=\operatorname{Tot}B_{\bullet,\bullet}^{d,N}(x),
 \qquad
 D_{\rm tot}:=d_{\rm bar}+(-1)^q\partial_\Gamma
 \quad\text{on }B_{q,r}^{d,N}.
\end{equation}
The functoriality of \(W\) and \(\Gamma\), together with
\(\partial_\Gamma^2=0\), gives
\begin{equation}\label{eq:pdfv2-coh-total-square-zero}
 D_{\rm tot}^2=0.
\end{equation}
In compact notation, but now already discharged by
\eqref{eq:pdfv2-coh-bar-bigraded}--
\eqref{eq:pdfv2-coh-bar-differential},
\begin{equation}\label{eq:pdfv2-coh-derived-coend}
 K_{d,N}(x)
 \simeq
 W_{d,N}\otimes^{\mathbf L}_{A_N[J_d(x)]}\Gamma_{d,N}.
\end{equation}

\section{Depth and terminal nonadic colimit}
\label{sec:pdfv2-coh-terminal}

The prefixes form the colimit.
\begin{equation}\label{eq:pdfv2-coh-k-infinity}
 K_{\infty,N}(x):=\varinjlim_d K_{d,N}(x).
\end{equation}
The nine-phase extension carries depth \(d\) to depth \(d+9\).
Only after forming \eqref{eq:pdfv2-coh-k-infinity} does it induce a
well-typed endomorphism
\begin{equation}\label{eq:pdfv2-coh-u-gamma9}
 U_{\Gamma_9,N}:K_{\infty,N}(x)
 \longrightarrow K_{\infty,N}(x).
\end{equation}
The terminal is
\begin{equation}\label{eq:pdfv2-coh-terminal-cone}
 \mathcal T_{{\rm coh},N}(x)
 :=\operatorname{Cone}\bigl(I-U_{\Gamma_9,N}\bigr).
\end{equation}
With the homological convention
\[
 \operatorname{Cone}(f)_q
 =K_{\infty,N,q-1}\oplus K_{\infty,N,q},
\]
its differential is
\begin{equation}\label{eq:pdfv2-coh-cone-differential}
 \partial_{\rm Cone}(a,b)
 =\bigl(-\partial_Ka,\partial_Kb+f(a)\bigr).
\end{equation}

Although the visible phase of \(h\) and \(\Gamma_9h\) coincides, the generators
\([h]\) and \([\Gamma_9h]\) are different because their ledgers and depths
are different. The terminal object preserves
\begin{equation}\label{eq:pdfv2-coh-terminal-difference}
 (I-U_{\Gamma_9,N})[h]=[h]-[\Gamma_9h].
\end{equation}
Precision reductions commute with the terminal:
\begin{equation}\label{eq:pdfv2-coh-u-reduction}
 r_{N+1,N}U_{\Gamma_9,N+1}
 =U_{\Gamma_9,N}r_{N+1,N}.
\end{equation}

\section{Terminal cohomological reader: incidence cycle and class in \(H_0(\operatorname{Cone}(I-U_{\Gamma_9,N}))\)}
\label{sec:pdfv2-coh-reader}

Each cell of \(x_{\leq d}\) carries a vector of incidence
\[
 \omega_{\nu,N}(x)
 \in A_N^{I_\nu\times J_\nu}=C_{\nu,N,0}
\]
and the complete terminal route \(\lambda_{\nu,x}\). The raw cycle is
\begin{equation}\label{eq:pdfv2-coh-z-cycle}
 z_{d,N}(x):=
 \sum_{\nu\in\operatorname{Cell}_d(x)}
 \epsilon_\nu[\lambda_{\nu,x}]
 \otimes\omega_{\nu,N}(x)
 \in B_{0,0}^{d,N}(x).
\end{equation}
As the total complex is homologically non-negative,
\begin{equation}\label{eq:pdfv2-coh-z-closed}
 D_{\rm tot}z_{d,N}(x)=0.
\end{equation}
Both the cycle and its class are preserved.
\[
 [z_{d,N}(x)]\in H_0(K_{d,N}(x)).
\]
After including it in \(K_{\infty,N}\), its terminal class is
\[
 [(0,z_{d,N}(x))]
 \in H_0(\mathcal T_{{\rm coh},N}(x)).
\]
The complete reader is
\begin{equation}\label{eq:pdfv2-coh-reader-full}
 \mathcal L_{{\rm coh},d,N}(x)
 :=\bigl(
 z_{d,N}(x),[z_{d,N}(x)],
 U_{\Gamma_9,N}z_{d,N}(x),
 [(0,z_{d,N}(x))]
 \bigr).
\end{equation}
The publication is not merely a quotient class: the representative, the route that produced it, its continuation and the terminal remain.

\section{Surviving fibres and corrections without a selector}
\label{sec:pdfv2-coh-surviving-fibers}

Let \(\mathscr S_\lambda^{\rm coh}\) be the free \(A_N\)-module generated by the
admissible finite histories at level \(\lambda=(d,N)\). The reader extends
linearly to
\begin{equation}\label{eq:pdfv2-coh-e-lambda}
 e_\lambda:\mathscr S_\lambda^{\rm coh}
 \longrightarrow\mathscr O_\lambda^{\rm coh},
 \qquad
 \mathscr O_\lambda^{\rm coh}
 :=H_0(K_{d,N})\oplus H_0(\mathcal T_{{\rm coh},N}).
\end{equation}
The surviving codomain is
\begin{equation}\label{eq:pdfv2-coh-surviving-output}
 \mathscr O_\infty^{\rm surv}
 :=\operatorname{Im}\!\left(
 e_\infty:\mathcal C_{\rm cont}^{\rm coh}
 \longrightarrow\varprojlim_\lambda\mathscr O_\lambda^{\rm coh}
 \right).
\end{equation}
For \(a=(a_\lambda)_\lambda\in\mathscr O_\infty^{\rm surv}\), define
\begin{equation}\label{eq:pdfv2-coh-fiber}
 \mathscr F_\lambda(a):=
 \left\{s_\lambda\in\mathscr S_\lambda^{\rm coh}:
 \begin{array}{l}
 e_\lambda(s_\lambda)=a_\lambda,\\
 \exists y\in\mathcal C_{\rm cont}^{\rm coh}\text{ with }
 y|_\lambda=s_\lambda\text{ and }e_\infty(y)=a
 \end{array}
 \right\}.
\end{equation}
These fibers are nonempty by the definition of the surviving codomain. If
\(\lambda'\geq\lambda\), the restriction is surjective:
\begin{equation}\label{eq:pdfv2-coh-fiber-surjection}
 \rho_{\lambda',\lambda}:\mathscr F_{\lambda'}(a)
 \twoheadrightarrow\mathscr F_\lambda(a).
\end{equation}
Indeed, a complete history that witnesses
\(s_\lambda\in\mathscr F_\lambda(a)\) provides by truncation an
antecedent at any later level. That antecedent is not chosen as
part of the object.

The complete multisect is
\begin{equation}\label{eq:pdfv2-coh-extension-multisection}
 \Sigma_{\lambda',\lambda}^{\rm coh}(s_\lambda;a)
 :=\{s_{\lambda'}\in\mathscr F_{\lambda'}(a):
 \rho_{\lambda',\lambda}s_{\lambda'}=s_\lambda\}.
\end{equation}
For a provisional prolongation \(\widetilde s_{\lambda'}\), let
\begin{equation}\label{eq:pdfv2-coh-discrepancy}
 \delta_{\lambda'}:=
 a_{\lambda'}-e_{\lambda'}(\widetilde s_{\lambda'}).
\end{equation}
The relation of all corrections is
\begin{equation}\label{eq:pdfv2-coh-correction-relation}
 \operatorname{Corr}_{\lambda',\lambda}
 (\widetilde s_{\lambda'},\delta_{\lambda'})
 :=\left\{\eta\in\mathscr S_{\lambda'}^{\rm coh}:
 \begin{array}{l}
 \rho_{\lambda',\lambda}\eta=0,\\
 e_{\lambda'}(\eta)=\delta_{\lambda'},\\
 \eta\text{ is supported on new cells},\\
 \eta\text{ preserves orientation, carry, memory and terminal}
 \end{array}
 \right\}.
\end{equation}
If two lifts are comparable, their difference
\(\eta=s_{\lambda'}-\widetilde s_{\lambda'}\) belongs to this relation
exactly when it satisfies the four preceding conditions. It is not asserted
that \eqref{eq:pdfv2-coh-correction-relation} is nonempty for an
arbitrary discrepancy: the governing non-emptiness is that of
\eqref{eq:pdfv2-coh-extension-multisection} on the surviving domain.

\iffalse
\section{The thirteen coordinates of the constraint}
\label{sec:pdfv2-coh-thirteen}

For \(x_\lambda\in\mathcal C_{{\rm cont},\lambda}^{\rm coh}\), let
\begin{equation}\label{eq:pdfv2-coh-d-package}
 \begin{aligned}
 D_{{\rm coh},\lambda}(x_\lambda):=\bigl(&
 \mathcal F_{{\rm coh},\lambda},
 \operatorname{Ext}_{\Gamma,{\rm coh},\lambda},
 \operatorname{Rel}_{{\rm coh},\lambda},
 \mathcal N_{{\rm coh},\lambda},
 \operatorname{or}_{{\rm coh},\lambda},
 \operatorname{Carr}_{{\rm coh},\lambda},
 \operatorname{Mem}_{{\rm coh},\lambda},\\
 &\partial_{{\rm coh},\lambda},
 \gamma_{{\rm coh},\lambda},
 \operatorname{Msect}_{{\rm coh},\lambda},
 \operatorname{Surv}_{{\rm coh},\lambda},
 \operatorname{Term}_{{\rm coh},\lambda},
 \mathcal L_{{\rm coh},\lambda}\bigr).
 \end{aligned}
\end{equation}
Their coordinates are:
\begin{enumerate}
 \item \(\mathcal F_{{\rm coh},\lambda}\): all the modulos
       \(A_N^{I_\nu}\), \(A_N^{J_\nu}\),
       \(A_N^{I_\nu\times J_\nu}\), the complejos \(C_{\nu,N}\), the
       modulos \(W_{d,N}(\nu)\) and \(\mathscr S_\lambda^{\rm coh}\).
 \item \(\operatorname{Ext}_{\Gamma,{\rm coh},\lambda}\): all maps
       \(\Gamma(u)\), coefficient reductions, truncations, extensions
       \(\Gamma_9\), and the relations \(\Sigma_{\lambda',\lambda}\).
 \item \(\operatorname{Rel}_{{\rm coh},\lambda}\):
       \(\kappa\), \(\delta\), the exact sequence, the bar faces, the
       coend relation, the orientation, and the composition laws.
 \item \(\mathcal N_{{\rm coh},\lambda}\):
       \[
       (d,N,3^N,a_\nu,b_\nu,(a_\nu-1)(b_\nu-1),
       \rho_{\Sigma,\nu},\rho_{\Pi,\nu},c_{\Sigma,\nu},c_{\Pi,\nu}).
       \]
 \item \(\operatorname{or}_{{\rm coh},\lambda}\): the signs
       \(\epsilon_\nu\), the port orderings, and the maps
       \(\Theta_{\nu,N}\).
 \item \(\operatorname{Carr}_{{\rm coh},\lambda}\): the matrices
       \(\beta_N\), the cocycles \(c_N(v,u)\), residues, and quotients,
       conserved separately.
 \item \(\operatorname{Mem}_{{\rm coh},\lambda}\): the inherited ledger and the
       ordered word of refinements, carries, corrections, and advances
       \(\Gamma_9\), not merely its aggregate value.
 \item \(\partial_{{\rm coh},\lambda}\):
       \(\kappa_{\nu,N}\), \(D_{\rm tot}\),
       \(\partial_{\rm Cone}\) and the boundaries of the cells.
 \item \(\gamma_{{\rm coh},\lambda}\): all composable chains
       \(\nu_0\to\cdots\to\nu_q\) and their terminal continuations.
 \item \(\operatorname{Msect}_{{\rm coh},\lambda}\): all cells,
       terminal routes, extensions \(\Sigma\), and corrections
       \(\operatorname{Corr}\), without a selector.
 \item \(\operatorname{Surv}_{{\rm coh},\lambda}\): the system of fibres
       \(\mathscr F_\lambda(a)\) and surjective restrictions.
 \item \(\operatorname{Term}_{{\rm coh},\lambda}\):
       \(\operatorname{Cone}(I-U_{\Gamma_9,N})\), together with \(z\), \(Uz\)
       and their difference.
 \item \(\mathcal L_{{\rm coh},\lambda}\):
       \((z,[z],Uz,[(0,z)])\) and all its reduction compatibilities.
\end{enumerate}

For \(\lambda'\geq\lambda\), the coordinated transport satisfies
\begin{equation}\label{eq:pdfv2-coh-d-naturality}
 u_{\lambda',\lambda}^{\rm coh}
 D_{{\rm coh},\lambda'}(x_{\lambda'})
 =D_{{\rm coh},\lambda}
  (\rho_{\lambda',\lambda}x_{\lambda'}).
\end{equation}

\section{Constraint graph and dependent object}
\label{sec:pdfv2-coh-graphs}

The total restriction is the graph
\begin{equation}\label{eq:pdfv2-coh-g-graph}
 \mathcal G_{\rm coh}:=\operatorname{Graph}(D_{\rm coh}),
 \qquad
 \operatorname{Res}_{\rm coh}(x):=(x,D_{\rm coh}x).
\end{equation}
Thus,
\begin{equation}\label{eq:pdfv2-coh-res-inverse}
 \pi_1\operatorname{Res}_{\rm coh}=I,
 \qquad
 \operatorname{Res}_{\rm coh}\pi_1=I_{\mathcal G_{\rm coh}}.
\end{equation}

The full extractor is
\begin{equation}\label{eq:pdfv2-coh-chi}
 \chi_{\rm coh}(x):=
 \bigl(
 \{C_{\nu,N}\},
 \{z_\lambda,[z_\lambda]\},
 \{\mathscr F_\lambda(e_\infty x)\},
 \{\Sigma_{\lambda',\lambda}\},
 \{\operatorname{Corr}_{\lambda',\lambda}\}
 \bigr).
\end{equation}
We define
\begin{equation}\label{eq:pdfv2-coh-xchi}
 X_{\chi,{\rm coh}}
 :=\{(\chi_{\rm coh}(x),x,D_{\rm coh}x):
 x\in\mathcal C_{\rm cont}^{\rm coh}\}
\end{equation}
and
\begin{equation}\label{eq:pdfv2-coh-prx}
 \operatorname{pr}_{X,{\rm coh}}(x,D_{\rm coh}x)
 :=(\chi_{\rm coh}(x),x,D_{\rm coh}x).
\end{equation}
The projection onto \((x,D_{\rm coh}x)\) is inverse on this object
dependent.

\section{Internal assembler and publisher}
\label{sec:pdfv2-coh-assembler-publisher}

The ambient space \(\mathscr M_{\rm coh}^{\rm amb}\) consists of systems
compatible of exact sequences of cells, categories, functors, modules of
path, totalizers, terminals, readers, fibers and multisections. The
assembler raw is
\begin{equation}\label{eq:pdfv2-coh-assembler}
 \begin{aligned}
 a_{\rm coh}:X_{\chi,{\rm coh}}&\longrightarrow
 \mathscr M_{\rm coh}^{\rm amb},\\
 a_{\rm coh}(\chi(x),x,Dx)
 &:={}
 \bigl(
 \{C_{\nu,N},\kappa_{\nu,N},\delta_{\nu,N}\},
 \{J_d,\Gamma,W,K\},\\
 &\hspace{31mm}
 \{U_{\Gamma_9,N},\mathcal T_N\},
 \{e_\lambda,\mathscr F_\lambda,\Sigma,\operatorname{Corr}\}
 \bigr)_x.
 \end{aligned}
\end{equation}
The macrostructure is another graph:
\begin{equation}\label{eq:pdfv2-coh-macro-graph}
 \mathfrak M_{\rm coh}:=\operatorname{Graph}(a_{\rm coh}),
 \qquad
 \mathcal A_{\rm coh}(z):=(z,a_{\rm coh}z).
\end{equation}

Let \(\mathscr Y_{\rm coh}^{\rm int}\) be the type of certificates containing
\begin{equation}\label{eq:pdfv2-coh-certificate-list}
 \begin{gathered}
 \text{la sucesión exacta \eqref{eq:pdfv2-coh-exact-sequence}},\\
 \kappa_{b,a}\tau=-P\kappa_{a,b},
 \qquad D_{\rm tot}^2=0,
 \qquad\partial_{\rm Cone}^2=0,\\
 rU_{\Gamma_9}=U_{\Gamma_9}r,
 \qquad
 e_\lambda\rho_{\lambda',\lambda}
 =r_{\lambda',\lambda}e_{\lambda'},\\
 \rho_{\lambda',\lambda}:\mathscr F_{\lambda'}(a)
 \twoheadrightarrow\mathscr F_\lambda(a),
 \qquad
 \Sigma_{\lambda',\lambda}(s_\lambda;a)\neq\varnothing,
 \end{gathered}
\end{equation}
together with the raw cycles, their classes and their terminals. The publisher
\begin{equation}\label{eq:pdfv2-coh-publisher}
 p_{\rm coh}^{\rm raw}:\mathfrak M_{\rm coh}
 \longrightarrow\mathscr Y_{\rm coh}^{\rm int}
\end{equation}
produces that certificate without eliminating the macrostructure. Finally,
\begin{equation}\label{eq:pdfv2-coh-output-graph}
 \mathcal C_{\rm coh}^{\rm HMT}
 :=\operatorname{Graph}(p_{\rm coh}^{\rm raw}),
 \qquad
 \mathcal P_{\rm coh}(m):=(m,p_{\rm coh}^{\rm raw}m).
\end{equation}

\section{Internal Theorem of Coherent Generation}
\label{sec:pdfv2-coh-theorem}

\begin{theorem}[Internal Generation of the Coherence Macrostructure]
\label{thm:pdfv2-coh-internal-generation}
On the explicit domain \(\mathcal C_{\rm cont}^{\rm coh}\), the chain
\[
 \mathcal C_{\rm cont}^{\rm coh}
 \xrightarrow{\operatorname{Res}_{\rm coh}}
 \mathcal G_{\rm coh}
 \xrightarrow{\operatorname{pr}_{X,{\rm coh}}}
 X_{\chi,{\rm coh}}
 \xrightarrow{\mathcal A_{\rm coh}}
 \mathfrak M_{\rm coh}
 \xrightarrow{\mathcal P_{\rm coh}}
 \mathcal C_{\rm coh}^{\rm HMT}
\]
naturally constructs the cells, their cones, the diagram of histories, the
two-sided bar, the totalization, the nonadic terminal, and the
surviving fibers. It conserves the thirteen coordinates and all corrections as
multisections. Each stage admits, on its image, a projection that recovers
the preceding stage.
\end{theorem}

\begin{proof}
The \cref{prop:pdfv2-coh-cell-exact} constructs each cell over \(A_N\), with
diagonal kernel and defect \((a_\nu-1)(b_\nu-1)\). The identity
\eqref{eq:pdfv2-coh-orientation} makes the oriented exchange an involutive map of
complexes. The formulas
\eqref{eq:pdfv2-coh-gamma-one}--\eqref{eq:pdfv2-coh-gamma-zero} prove that
each refinement is a map of complexes, whereas
\eqref{eq:pdfv2-coh-kappa-reduction} proves compatibility in precision.

Precomposition of paths makes \(W\). The identities
functorial of \(W\) and \(\Gamma\) imply
\(d_{\rm bar}^2=0\); its commutation with the cellular differential gives
\(D_{\rm tot}^2=0\). Thus, \(K_{d,N}\) is constructed by the bar formulae,
not by to nominal invocation of the coend.

The extension \(\Gamma_9\) increases the depth. Therefore it induces the
endomorphism \(U_{\Gamma_9,N}\) only after passage to the colimit
\(K_{\infty,N}\). The cone
\(\operatorname{Cone}(I-U_{\Gamma_9,N})\) then conserves the difference
\([h]-[\Gamma_9h]\) and does not identify phase repetition with repetition
of state.

The cycle \(z_{d,N}(x)\) is formed with all the cells, signs, routes, and incidences of the history. Truncation and reduction commute with the reader. If \(s_\lambda\in\mathscr F_\lambda(a)\), the complete history appearing in the very definition of the fiber provides an antecedent at every level \(\lambda'\geq\lambda\); this proves the surjectivity \eqref{eq:pdfv2-coh-fiber-surjection}. The object preserves the entire set \(\Sigma_{\lambda',\lambda}\), not a choice. Differences between lifts are registered through \eqref{eq:pdfv2-coh-correction-relation}, with support, orientation, carry, memory, and terminal made explicit.

Finally, \(\mathcal G_{\rm coh}\), \(\mathfrak M_{\rm coh}\), and
\(\mathcal C_{\rm coh}^{\rm HMT}\) are graphs, whereas
\(X_{\chi,{\rm coh}}\) retains \(x\) and \(D_{\rm coh}x\) as coordinates.
The first projection of each object recovers its predecessor. Therefore, the
complete composition does not replace history with class, multisection with
selector, or terminal object with visible phase.
\end{proof}

\section{Probative origin and forgers}
\label{sec:pdfv2-coh-provenance-falsifiers}

\paragraph{Result recovered.}
Belonging to the already constructed architecture are the cell \(\kappa_{a,b}\), its
diagonal kernel and defect, the orientation
\(\kappa_{b,a}\tau=-P\kappa_{a,b}\), the APP carry, the cone of the cell, the
category of histories, the totalization by coend, the action \(\Gamma_9\), the
terminal \(\operatorname{Cone}(I-U_{\Gamma_9})\), and the architecture of
surviving extension.

\paragraph{New Formalization.}
Explicit formalizations include the exact sequence \eqref{eq:pdfv2-coh-exact-sequence} over \(A_N\), the separation between depth and precision, the complete faces of the bar, formation of the terminal object only after the colimit, packaging of the thirteen coordinates, corrections as a relation without a selector, and the four dependent objects or graphs that prevent loss of genealogy.

\begin{criterion}[Internal falsifiers of the branch \(\mathrm{coh}\)]
\label{crit:pdfv2-coh-falsifiers}
The construction fails if any of the following occurs:
\begin{enumerate}
 \item the sequence is not exact
       \eqref{eq:pdfv2-coh-exact-sequence};
 \item \(\kappa_{b,a}\tau=-P\kappa_{a,b}\) fails;
 \item some refinement does not satisfy
       \(\kappa_{\nu'}\Gamma(u)_1=\Gamma(u)_0\kappa_\nu\);
 \item the composition carry does not satisfy
       \eqref{eq:pdfv2-coh-carry-cocycle};
 \item \(D_{\rm tot}^2\neq0\);
 \item reduction, reader, or terminal do not commute with \(U_{\Gamma_9}\);
 \item \(\operatorname{Cone}(I-U_{\Gamma_9})\) is formed by previously identifying
       the depths \(d\) and \(d+9\);
 \item a fiber declared to be surviving is empty, or a restriction is not
       surjective;
 \item is replaced \(\Sigma_{\lambda',\lambda}\) by a single
       lift;
 \item only \([z]\) is published, and \(z\), route, carry, memory, or
       terminal is eliminated;
 \item any of the thirteen coordinates is omitted.
\end{enumerate}
In the last case, the binding diagnosis is: \emph{object replaced;
conclusion not transferable}.
\end{criterion}

\begin{warningbox}
This chapter concludes only the internal construction of
\(\mathcal C_{\rm coh}^{\rm HMT}\). It introduces neither a classical space,
nor a classical class, nor a classical publication. Any later intertwiner
belongs to another document and does not govern the existence of this
macrostructure.
\end{warningbox}
\fi
